\section{The Hermes Planner}
\label{sec-framework}

\hermes has four components: an entity linker, an explorer that applies
actions and observes their effect, a verifier that decides validity and
goal satisfaction, and a planner that chooses actions.
Algorithm~\ref{alg-plan} in Appendix~B gives the loop.

\stitle{Actions}.
The action space is the algebra of Table~\ref{tab-ops}, plus three
kinds of auxiliary actions.
\emph{Graph reads} leave the plan unchanged: \kw{lookup} returns
entities whose label contains a string, and \kw{inspect} lists the
predicates around an entity or used by the instances of a class, with
counts and an example value.
\emph{Control} actions are \kw{goal}, which states $\sigma$;
\kw{undo}, which restores the previous plan; and \kw{finish}.
Each action is a JSON object, so a planner emits one action per turn
and the explorer parses it into a typed command.

\stitle{Observations}.
After a plan-changing action, the explorer renders the new plan, checks
that it parses, asks the verifier for validity, and only then executes
two queries on the graph store: the number of solutions, as a
\kw{COUNT} over the plan used as a subquery, and a sample of rows under
a \kw{LIMIT} (five by default).
The observation reports the count, the previous count when it changed,
the sample, and any error.
An action whose plan does not parse, is not valid, or fails to run is
not committed.
\kw{preview} returns the same observation without committing, the
analogue of costing a rewrite before choosing it.
Executions are memoized, so revisiting a plan costs nothing.

\stitle{Verifier}.
The verifier computes $\Gamma_\pi$ from the schema, checks the
conditions of Section~\ref{sec-problem}, and returns errors, warnings,
the typing context, and the conditions of $\sigma$ that $\pi$ does not
yet satisfy (the \emph{open goals}).
Error messages name the offending term and, for an unknown class or
property, suggest the closest term in the schema.
Algorithm~\ref{alg-verify} in Appendix~B gives the procedure.

\stitle{Planner}.
The planner is an LLM prompted, at every turn, with the action
catalogue, the schema vocabulary, the entities whose labels occur in
the question, the history of actions and observations, and the current
plan as SPARQL.
With the verifier enabled, the prompt also shows the typing context and
the open goals, and \kw{finish} is refused while an error or an open
goal remains.
An action identical to the previous one is not executed.
The loop ends on \kw{finish} or after a step budget (15), and returns
the last committed plan that runs.

\stitle{Other policies in the same space}.
Existing generators are policies over these actions.
One-shot generation is a single macro action that sets every field of
the plan at once.
Constraint-guided search~\cite{wang2026uniqgen} is a beam search over
\kw{remove} and \kw{triple}: a Chase phase removes constraints until the
plan covers a reference answer set, and a Backchase phase adds them
back until it returns only reference answers~\cite{deutsch1999chase}.
Placing them in one space lets the evaluation hold the graph store,
the scorer, and the model fixed while varying only the policy.

\stitle{Implementation}.
\hermes is written in Python.
Graph stores share one interface that returns SPARQL 1.1 JSON results,
so the same planner runs on local RDF files loaded into
Oxigraph~\cite{oxigraph} or rdflib, or on a remote SPARQL endpoint.
Model completions are cached on disk, so a rerun replays the same
outputs, and each run records the model, the number of model calls and
graph executions, and the commit it started from.
